\documentclass[12pt]{article}

\usepackage[margin=1in]{geometry}
\usepackage{amsmath,amsthm,amssymb}
\usepackage{mathtools}
\usepackage{hyperref}
\usepackage{enumitem}
\usepackage{pgfplots}
\pgfplotsset{compat=1.18}

% Theorem environments
\newtheorem{theorem}{Theorem}[section]
\newtheorem{proposition}[theorem]{Proposition}
\newtheorem{lemma}[theorem]{Lemma}
\newtheorem{corollary}[theorem]{Corollary}
\theoremstyle{definition}
\newtheorem{definition}[theorem]{Definition}
\newtheorem{remark}[theorem]{Remark}
\newtheorem{example}[theorem]{Example}

% Shortcuts
\newcommand{\R}{\mathbb{R}}
\newcommand{\E}{\mathbb{E}}
\newcommand{\tr}{\operatorname{tr}}
\newcommand{\Id}{I_d}
\newcommand{\Spp}{S_d^{++}}
\newcommand{\Sp}{S_d^{+}}

\title{The Gaussian Schr\"odinger--Bass Bridge:\\Explicit Optimizers and the Solvability Boundary}
\author{}
\date{}

\begin{document}
\maketitle
\thispagestyle{empty}

\begin{abstract}
We provide a self-contained derivation, via Pontryagin's maximum principle,
of the explicit Gaussian solution of the Schr\"odinger--Bass Bridge (SBB)
problem introduced by Henry-Labord\`ere, Loeper, Mazhar, Pham, and
Touzi~\cite{HLMPT26} and studied independently in~\cite{PDE26,LSBB26}.
The SBB interpolates between the classical Schr\"odinger bridge
($\beta\to\infty$) and the Bass martingale transport ($\beta\to 0$) via
a parameter $\beta>0$.
Our Pontryagin approach yields the adjoint equation $\dot p = 2p^2$
in closed form, reducing the boundary value problem to a quadratic equation.
We identify the \emph{solvability boundary}: the critical target variance
$\gamma_T^* = (\gamma_0 + \beta T^2)/(1+2\beta T)^2$
below which the Gaussian SBB admits no finite-cost solution.
Going beyond the Gaussian setting, we prove that $\gamma_T^*$ is
a \emph{second-moment lower bound} for solvability of the SBB with
arbitrary (non-Gaussian) centered marginals.
We characterize the degenerate boundary regime, where the optimal
diffusion vanishes terminally and the value function has a jump
discontinuity.
A numerical example with non-commuting covariance matrices demonstrates
anisotropic costate coupling in dimensions $d \ge 2$.
\end{abstract}

\setcounter{page}{1}

%%%%%%%%%%%%%%%%%%%%%%%%%%%%%%%%%%%%%%%%%%%%%%%%%%%%%%%%%%%%%
\section{Introduction}\label{sec:intro}
%%%%%%%%%%%%%%%%%%%%%%%%%%%%%%%%%%%%%%%%%%%%%%%%%%%%%%%%%%%%%

The Schr\"odinger bridge problem, originating in Schr\"odinger's 1931 thought
experiment~\cite{CGP16}, seeks the most likely evolution of a cloud of
independent Brownian particles conditioned to match prescribed initial and
terminal distributions.
In modern formulation, given marginals $\mu_0, \mu_T \in \mathcal{P}_2(\R^d)$,
one minimizes the relative entropy of a semimartingale law with respect to
Wiener measure, subject to the marginal constraints.
The Bass construction~\cite{BHK20} provides a different paradigm: one seeks a
martingale interpolation between two marginals by controlling only the
diffusion coefficient of a local martingale.

Henry-Labord\`ere, Loeper, Mazhar, Pham, and Touzi~\cite{HLMPT26} introduced
the \emph{Schr\"odinger--Bass Bridge} (SBB) problem, which interpolates between
these two frameworks via a single parameter $\beta > 0$.
The SBB simultaneously controls both the drift $\alpha_t$ and the diffusion
coefficient $a_t = \sigma_t \sigma_t^\top$ of a semimartingale, penalizing
drift via $\frac{1}{2}|\alpha_t|^2$ and diffusion deviations from the identity
via $\frac{\beta}{2}\|a_t - I_d\|_F^2$.
As $\beta\to\infty$, the diffusion coefficient is forced to $I_d$ and one
recovers the classical Schr\"odinger bridge; as $\beta\to 0$, the drift penalty
dominates and one approaches the Bass martingale transport.

Independently, Backhoff-Veraguas, Beiglb\"ock, Bifronte, and Ley~\cite{BBBL26}
studied the martingale Schr\"odinger bridge from a related perspective, while
Hasenbichler, Pammer, and Thonhauser~\cite{HPT26} developed a weak transport
approach.
A PDE derivation of the one-dimensional SBB system, including explicit
solutions via Legendre transforms and the heat equation, was given
in~\cite{PDE26}; computational aspects were addressed by Alouadi
et~al.~\cite{LSBB26}.
The Gaussian Schr\"odinger bridge (without Bass interpolation) has a
well-known closed form~\cite{Bunne23,CGP16};
generalizations of the Bass martingale to arbitrary reference measures
were developed by Tschiderer~\cite{Tsch24}.

A key finding of~\cite{HLMPT26} is that the SBB is well-posed under
$\beta T > 1$.
The Gaussian case was solved explicitly in~\cite{PDE26} via PDE methods.
However, the following questions remained open:
the precise solvability boundary for Gaussian marginals (and its
relationship to the $\beta T > 1$ condition), whether this boundary
extends to non-Gaussian marginals, the structure of the degenerate
regime at the boundary, and the behavior of non-commuting covariance
matrices in $d \ge 2$ dimensions.

\medskip

\noindent\textbf{Main contributions.}
We provide a self-contained derivation of the Gaussian SBB via Pontryagin's
maximum principle (complementing the PDE approach of~\cite{PDE26}) and
address the open questions above.
Our results are organized as follows:

\begin{enumerate}[label=(\roman*)]
\item \textbf{Theorem~\ref{thm:1d}} gives the explicit optimizer for the
one-dimensional Gaussian SBB.
The optimal drift and volatility are expressed in closed form through a
scalar parameter $p_0$ determined by a quadratic equation.

\item \textbf{Theorem~\ref{thm:solvability}} identifies the
\emph{solvability boundary}: the critical target variance
$\gamma_T^* = (\gamma_0 + \beta T^2)/(1 + 2\beta T)^2$,
below which no Gaussian SBB solution exists.
The Gaussian SBB is solvable if and only if
$\gamma_T \ge \gamma_T^*$.
This threshold clarifies the role of the condition
$\beta T > 1$ from~\cite{HLMPT26}: it ensures universal solvability
across all marginals, but specific marginal pairs may be solvable
even when $\beta T \le 1$.

\item \textbf{Theorem~\ref{thm:multid}} extends the results to $d$
dimensions, where the adjoint equation becomes the matrix Riccati
equation $\dot P = 2P^2$ with solution
$P(t) = P_0(I - 2P_0 t)^{-1}$.
When $\Gamma_0$ and $\Gamma_T$ commute, the $d$-dimensional problem
decomposes into $d$ independent scalar problems.

\item \textbf{Proposition~\ref{prop:limits}} characterizes the limiting
regimes $\beta\to\infty$ (Schr\"odinger bridge) and $\beta\to 0$
(Bass martingale transport).

\item \textbf{Theorem~\ref{thm:nongaussian}} establishes the Gaussian
solvability boundary as a \emph{universal lower bound}: for any
centered marginals $\mu_0, \mu_T$ (not necessarily Gaussian) with
second moment $\gamma_T < \gamma_T^*$, the SBB cost is infinite.

\item \textbf{Proposition~\ref{prop:boundary}} characterizes the
\emph{degenerate boundary regime}: at $\gamma_T = \gamma_T^*$,
the optimal cost is finite but the terminal diffusion vanishes,
with a jump discontinuity in the value function.

\item \textbf{Example~\ref{ex:noncommuting}} demonstrates the
non-commuting multi-dimensional case numerically, revealing
anisotropic costate structure inaccessible to the commuting
decomposition.
\end{enumerate}

\medskip

\noindent\textbf{Notation.}
We write $S_d$ for the space of $d\times d$ real symmetric matrices,
$\Spp$ for positive definite matrices, and $\Sp$ for positive semidefinite
matrices.
For $M \in S_d$, we write $M > 0$ (resp.\ $M \ge 0$) for $M \in \Spp$
(resp.\ $M \in \Sp$).
The trace is $\tr(M)$.
For a function $f(t)$, we write $\dot f$ for $df/dt$.
We use $|\cdot|$ for the Euclidean norm and $\|\cdot\|_F$ for the Frobenius
norm.

%%%%%%%%%%%%%%%%%%%%%%%%%%%%%%%%%%%%%%%%%%%%%%%%%%%%%%%%%%%%%
\section{Setup and Gaussian Reduction}\label{sec:setup}
%%%%%%%%%%%%%%%%%%%%%%%%%%%%%%%%%%%%%%%%%%%%%%%%%%%%%%%%%%%%%

\subsection{The Schr\"odinger--Bass Bridge Problem}

Let $\mu_0, \mu_T \in \mathcal{P}_2(\R^d)$ be probability measures with
finite second moments.
Consider a filtered probability space
$(\Omega, \mathcal{F}, (\mathcal{F}_t)_{t\in[0,T]}, \mathbb{P})$ supporting
a $d$-dimensional Brownian motion $(W_t)_{t\in[0,T]}$.

\begin{definition}[SBB Problem]\label{def:sbb}
The Schr\"odinger--Bass Bridge problem is
\begin{equation}\label{eq:sbb}
\inf_{(\alpha, \sigma)} \int_0^T \E\!\left[\frac{1}{2}|\alpha_t|^2 +
  \frac{\beta}{2}\|\sigma_t \sigma_t^\top - I_d\|_F^2\right] dt,
\end{equation}
subject to the dynamics
\begin{equation}\label{eq:dynamics}
dX_t = \alpha_t \, dt + \sigma_t \, dW_t, \qquad
X_0 \sim \mu_0, \quad X_T \sim \mu_T,
\end{equation}
where $\alpha_t$ is an $\R^d$-valued progressively measurable drift
and $\sigma_t$ is an $\R^{d\times d}$-valued progressively measurable
diffusion coefficient.
\end{definition}

\subsection{Gaussian Reduction}

When $\mu_0 = \mathcal{N}(m_0, \Gamma_0)$ and
$\mu_T = \mathcal{N}(m_T, \Gamma_T)$ are Gaussian, we may restrict to
Gaussian processes without loss of optimality.

\begin{proposition}[Gaussian Ansatz]\label{prop:gaussian}
Among all semimartingales satisfying~\eqref{eq:dynamics} with Gaussian
marginals, it suffices to consider processes with deterministic drift
and diffusion coefficients of the form $\alpha_t = A_t X_t + b_t$ and
$\sigma_t = \Sigma_t$ (deterministic).
The mean $m_t = \E[X_t]$ and covariance
$\Gamma_t = \E[(X_t - m_t)(X_t - m_t)^\top]$ then satisfy
\begin{equation}\label{eq:mean-cov}
\dot m = A_t m_t + b_t, \qquad
\dot\Gamma = A_t \Gamma_t + \Gamma_t A_t^\top + \Sigma_t \Sigma_t^\top.
\end{equation}
The cost~\eqref{eq:sbb} separates into a mean component and a covariance
component; the mean component is minimized by standard Schr\"odinger bridge
arguments, and the covariance component becomes
\begin{equation}\label{eq:cov-cost}
J = \int_0^T \left[\frac{1}{2}\tr(A_t^2 \Gamma_t) +
  \frac{\beta}{2}\|\Sigma_t \Sigma_t^\top - I_d\|_F^2\right] dt.
\end{equation}
\end{proposition}

\begin{proof}
The drift $\alpha_t = A_t X_t + b_t$ gives
$\E[|\alpha_t|^2] = \tr(A_t \Gamma_t A_t^\top) + |A_t m_t + b_t|^2$.
By the strict convexity of $(\alpha, a) \mapsto \frac{1}{2}|\alpha|^2
+ \frac{\beta}{2}\|a - I\|_F^2$ and Jensen's inequality applied to the
conditional expectations
$E[|\alpha_t|^2 \mid \Gamma_t] \ge |E[\alpha_t \mid \Gamma_t]|^2$ and
$E[(\sigma_t^2 - 1)^2] \ge (E[\sigma_t^2] - 1)^2$,
the cost is minimized by deterministic (state-independent) controls,
which corresponds to the Gaussian ansatz
(cf.~\cite{CGP16} for the analogous argument in the Schr\"odinger bridge).
The mean dynamics decouple from the covariance dynamics, and the mean
optimization reduces to a standard Schr\"odinger bridge~\cite{CGP16},
solved by $b_t = (m_T - m_0)/T - A_t m_t$.
The remaining cost depends only on $(A_t, \Sigma_t)$ and $\Gamma_t$ as
stated.
\end{proof}

For the remainder of the paper, we focus on the covariance optimization
problem~\eqref{eq:cov-cost} subject to the covariance
dynamics~\eqref{eq:mean-cov}, assuming without loss of generality that
$m_0 = m_T = 0$.

%%%%%%%%%%%%%%%%%%%%%%%%%%%%%%%%%%%%%%%%%%%%%%%%%%%%%%%%%%%%%
\section{One-Dimensional Gaussian SBB}\label{sec:1d}
%%%%%%%%%%%%%%%%%%%%%%%%%%%%%%%%%%%%%%%%%%%%%%%%%%%%%%%%%%%%%

In one dimension, we write $\gamma_t = \E[X_t^2]$ for the variance (with
zero mean).
The dynamics and cost become
\begin{equation}\label{eq:1d-dynamics}
\dot\gamma = 2\alpha_t \gamma_t + \sigma_t^2, \qquad
J = \int_0^T \left[\frac{1}{2}\alpha_t^2 \gamma_t +
  \frac{\beta}{2}(\sigma_t^2 - 1)^2\right] dt,
\end{equation}
with boundary conditions $\gamma_0$ and $\gamma_T$ prescribed.

\begin{remark}
The Gaussian SBB optimizer was first derived in~\cite{PDE26}
via PDE methods (Legendre transforms and the heat equation).
The following theorem provides an alternative derivation via
Pontryagin's maximum principle, which yields the same solution
in a form particularly suited to the solvability analysis of
Sections~\ref{sec:solvability}--\ref{sec:beyond}.
\end{remark}

\begin{theorem}[One-Dimensional Optimizer]\label{thm:1d}
Let $\gamma_0, \gamma_T > 0$, $\beta > 0$, and $T > 0$.
Define
\begin{equation}\label{eq:quadratic}
v = \frac{T + \sqrt{T^2 + 4\bigl(\gamma_T + \tfrac{1}{4\beta}\bigr)
  \bigl(\gamma_0 + \tfrac{1}{4\beta}\bigr)}}{2\bigl(\gamma_T +
  \tfrac{1}{4\beta}\bigr)},
\end{equation}
and set $p_0 = (v-1)/(2vT)$.
Then the unique optimal controls are
\begin{equation}\label{eq:opt-controls}
\alpha^*(t) = \frac{-2p_0}{1 - 2p_0 t}, \qquad
\sigma^{*2}(t) = 1 - \frac{p_0}{\beta(1 - 2p_0 t)}.
\end{equation}
The optimal costate is $p(t) = p_0/(1 - 2p_0 t)$, and the
optimal variance trajectory is determined by $\dot\gamma =
2\alpha^* \gamma + \sigma^{*2}$ with $\gamma(0) = \gamma_0$.
The conserved Hamiltonian along the optimal trajectory is
\begin{equation}\label{eq:hamiltonian-value}
H^* = -2p_0^2 \gamma_0 + p_0 - \frac{p_0^2}{2\beta}.
\end{equation}
\end{theorem}

\begin{proof}
We apply Pontryagin's maximum principle to the optimal control
problem~\eqref{eq:1d-dynamics}.

\medskip\noindent
\textbf{Step 1: Hamiltonian and first-order conditions.}
The control Hamiltonian is
\begin{equation}\label{eq:hamiltonian}
H = \frac{1}{2}\alpha^2 \gamma + \frac{\beta}{2}(\sigma^2 - 1)^2
  + p(2\alpha\gamma + \sigma^2),
\end{equation}
where $p$ is the costate variable conjugate to $\gamma$.
The first-order optimality conditions $\partial H/\partial\alpha = 0$
and $\partial H/\partial(\sigma^2) = 0$ give
\begin{equation}\label{eq:foc}
\alpha\gamma + 2p\gamma = 0 \implies \alpha^* = -2p, \qquad
\beta(\sigma^2 - 1) + p = 0 \implies \sigma^{*2} = 1 - \frac{p}{\beta}.
\end{equation}
The condition $\sigma^{*2} \ge 0$ requires $p \le \beta$.

\medskip\noindent
\textbf{Step 2: Adjoint equation.}
The adjoint equation is
\[
\dot p = -\frac{\partial H}{\partial \gamma}
  = -\frac{1}{2}\alpha^{*2} - 2p\alpha^*
  = -\frac{1}{2}(4p^2) - 2p(-2p)
  = -2p^2 + 4p^2 = 2p^2.
\]
Hence
\begin{equation}\label{eq:adjoint}
\dot p = 2p^2.
\end{equation}
This is a Bernoulli equation with the explicit solution
\begin{equation}\label{eq:p-sol}
p(t) = \frac{p_0}{1 - 2p_0 t},
\end{equation}
where $p_0 = p(0)$.

\medskip\noindent
\textbf{Step 3: Substituting back.}
With $u(t) = 1 - 2p_0 t$, the optimal controls
become~\eqref{eq:opt-controls}.
The variance dynamics under the optimal controls are
\begin{equation}\label{eq:gamma-opt}
\dot\gamma = \frac{-4p_0}{u} \gamma + 1 - \frac{p_0}{\beta u}.
\end{equation}

\medskip\noindent
\textbf{Step 4: Integrating factor.}
Define $\varphi(t) = \gamma(t)/u(t)^2$.
Using $\dot u = -2p_0$, we compute
\[
\dot\varphi = \frac{\dot\gamma}{u^2} - \frac{2\gamma \dot u}{u^3}
  = \frac{1}{u^2}\!\left(\dot\gamma + \frac{4p_0 \gamma}{u}\right)
  = \frac{1}{u^2}\!\left(1 - \frac{p_0}{\beta u}\right)
  = \frac{1}{u^2} - \frac{p_0}{\beta u^3}.
\]
Integrating from $0$ to $T$, with $u_T = u(T) = 1 - 2p_0 T$:
\begin{align}
\frac{\gamma_T}{u_T^2} - \gamma_0
  &= \int_0^T \!\left(\frac{1}{u^2} - \frac{p_0}{\beta u^3}\right) dt
  \notag\\
  &= \frac{1}{2p_0}\!\left(\frac{1}{u_T} - 1\right)
    - \frac{1}{4\beta}\!\left(\frac{1}{u_T^2} - 1\right),
    \label{eq:integrated}
\end{align}
where we used $dt = -du/(2p_0)$ and the elementary antiderivatives
$\int u^{-2}\,du = -u^{-1}$ and $\int u^{-3}\,du = -u^{-2}/2$.

\medskip\noindent
\textbf{Step 5: Boundary value problem.}
Setting $v = 1/u_T = 1/(1 - 2p_0 T)$, equation~\eqref{eq:integrated}
becomes
\[
\gamma_T v^2 - \gamma_0 = \frac{v - 1}{2p_0} -
  \frac{v^2 - 1}{4\beta}.
\]
Since $v = 1/(1 - 2p_0 T)$, we have $2p_0 = (v-1)/(vT)$, so
$1/(2p_0) = vT/(v-1)$.
From~\eqref{eq:integrated}, multiplying through by $u_T^2$ and
rearranging with $v = 1/u_T$:
\begin{align*}
\gamma_T - \gamma_0 u_T^2
  &= \frac{u_T^2}{2p_0}\left(\frac{1}{u_T} - 1\right)
    - \frac{u_T^2}{4\beta}\left(\frac{1}{u_T^2} - 1\right) \\
  &= \frac{u_T(1 - u_T)}{2p_0} - \frac{1 - u_T^2}{4\beta}.
\end{align*}
Since $1 - u_T = 2p_0 T$, the first term becomes
$u_T \cdot T = T/v$.
For the second term, $1 - u_T^2 = (1 - u_T)(1 + u_T) =
2p_0 T(1 + 1/v) = 2p_0 T(v+1)/v$.
Thus
\[
\gamma_T - \frac{\gamma_0}{v^2} = \frac{T}{v}
  - \frac{2p_0 T(v+1)}{4\beta v}
  = \frac{T}{v} - \frac{(v-1)(v+1)}{4\beta v^2},
\]
where we used $2p_0 T = 1 - 1/v = (v-1)/v$.
Multiplying by $v^2$:
\[
\gamma_T v^2 - \gamma_0 = Tv - \frac{v^2 - 1}{4\beta}.
\]
Rearranging:
\begin{equation}\label{eq:bvp}
v^2\!\left(\gamma_T + \frac{1}{4\beta}\right) - Tv
  - \left(\gamma_0 + \frac{1}{4\beta}\right) = 0.
\end{equation}
This is a quadratic in $v$ with coefficients
$a = \gamma_T + 1/(4\beta) > 0$, $b = -T$,
$c = -(\gamma_0 + 1/(4\beta)) < 0$.
The discriminant is
$T^2 + 4(\gamma_T + 1/(4\beta))(\gamma_0 + 1/(4\beta)) > 0$,
so there is a unique positive root given by~\eqref{eq:quadratic}.

\medskip\noindent
\textbf{Step 6: Recovery of $p_0$.}
From $v = 1/(1 - 2p_0 T)$, we obtain
\begin{equation}\label{eq:p0}
p_0 = \frac{v - 1}{2vT}.
\end{equation}
Since $v > 0$ is determined by $(\gamma_0, \gamma_T, \beta, T)$
via~\eqref{eq:quadratic}, all optimal controls are explicit.

\medskip\noindent
\textbf{Step 7: Hamiltonian conservation.}
Substituting $\alpha^* = -2p_0$, $\sigma^{*2} = 1 - p_0/\beta$,
and $\gamma(0) = \gamma_0$ into~\eqref{eq:hamiltonian} at $t = 0$:
\begin{align*}
H^* &= \frac{1}{2}(4p_0^2)\gamma_0 + \frac{\beta}{2}
  \left(\frac{p_0}{\beta}\right)^2 + p_0(2(-2p_0)\gamma_0 +
  1 - p_0/\beta) \\
  &= 2p_0^2 \gamma_0 + \frac{p_0^2}{2\beta} - 4p_0^2\gamma_0
    + p_0 - \frac{p_0^2}{\beta} \\
  &= -2p_0^2\gamma_0 + p_0 - \frac{p_0^2}{2\beta}.
\end{align*}
Since the Hamiltonian has no explicit time dependence (through the
state-costate formulation, $dH/dt = \partial H/\partial t = 0$),
$H^*$ is constant along the optimal trajectory.
\end{proof}

%%%%%%%%%%%%%%%%%%%%%%%%%%%%%%%%%%%%%%%%%%%%%%%%%%%%%%%%%%%%%
\section{Solvability Boundary}\label{sec:solvability}
%%%%%%%%%%%%%%%%%%%%%%%%%%%%%%%%%%%%%%%%%%%%%%%%%%%%%%%%%%%%%

The admissibility constraint $\sigma^{*2}(t) \ge 0$ for all
$t \in [0,T]$ imposes restrictions on which boundary data
$(\gamma_0, \gamma_T)$ admit a solution.

\begin{theorem}[Solvability Boundary]\label{thm:solvability}
Let $\beta > 0$, $T > 0$, and $\gamma_0 > 0$.
Define the critical target variance
\begin{equation}\label{eq:critical}
\gamma_T^* = \frac{\gamma_0 + \beta T^2}{(1 + 2\beta T)^2}.
\end{equation}
Then:
\begin{enumerate}[label=(\alph*)]
\item The Gaussian SBB with data $(\gamma_0, \gamma_T, \beta, T)$
  is solvable if and only if $\gamma_T \ge \gamma_T^*$.

\item At $\gamma_T = \gamma_T^*$, the optimal terminal volatility
  degenerates: $\sigma^*(T) = 0$.

\item $\gamma_T^* > 0$ for all $\beta, T > 0$.

\item As $\beta \to \infty$: $\gamma_T^* \to 0$.

\item As $\beta \to 0$: $\gamma_T^* \to \gamma_0$.
\end{enumerate}
\end{theorem}

\begin{proof}
\textbf{Step 1: Admissibility constraint.}
The optimal volatility is $\sigma^{*2}(t) = 1 - p_0/(\beta u(t))$
where $u(t) = 1 - 2p_0 t$.
Since $u$ is monotone and $u(T) = 1/v$ is its most extreme value
(for $p_0 > 0$, $u$ is decreasing), the binding constraint is at $t = T$:
\[
\sigma^{*2}(T) = 1 - \frac{p_0 v}{\beta} = 1 - \frac{v - 1}{2\beta T v}
  \cdot v = 1 - \frac{v - 1}{2\beta T} \ge 0.
\]
This gives the admissibility bound
\begin{equation}\label{eq:admissibility}
v \le 1 + 2\beta T.
\end{equation}

\medskip\noindent
\textbf{Step 2: Critical threshold.}
The critical case is $v = v^* := 1 + 2\beta T$.
Substituting into the BVP~\eqref{eq:bvp}:
\[
(1 + 2\beta T)^2 \!\left(\gamma_T^* + \frac{1}{4\beta}\right)
  - T(1 + 2\beta T) - \left(\gamma_0 + \frac{1}{4\beta}\right) = 0.
\]
Solving for $\gamma_T^*$:
\begin{align*}
\gamma_T^* &= \frac{T(1 + 2\beta T) + \gamma_0 + \frac{1}{4\beta}}
  {(1 + 2\beta T)^2} - \frac{1}{4\beta} \\
  &= \frac{T(1 + 2\beta T) + \gamma_0 + \frac{1}{4\beta}
    - \frac{(1 + 2\beta T)^2}{4\beta}}{(1 + 2\beta T)^2}.
\end{align*}
We expand the numerator:
\begin{align*}
&T(1 + 2\beta T) + \gamma_0 + \frac{1}{4\beta}
  - \frac{1 + 4\beta T + 4\beta^2 T^2}{4\beta} \\
&= T + 2\beta T^2 + \gamma_0 + \frac{1}{4\beta}
  - \frac{1}{4\beta} - T - \beta T^2 \\
&= \gamma_0 + \beta T^2.
\end{align*}
Therefore
\[
\gamma_T^* = \frac{\gamma_0 + \beta T^2}{(1 + 2\beta T)^2}.
\]

\medskip\noindent
\textbf{Step 3: Solvability direction.}
We show $v$ is strictly decreasing in $\gamma_T$.
By the implicit function theorem applied to
$F(v, \gamma_T) := (\gamma_T + \tfrac{1}{4\beta})v^2 - Tv
- (\gamma_0 + \tfrac{1}{4\beta}) = 0$:
\[
\frac{dv}{d\gamma_T} = -\frac{\partial F/\partial \gamma_T}
{\partial F/\partial v}
= -\frac{v^2}{2(\gamma_T + \tfrac{1}{4\beta})v - T}.
\]
Since $v$ is the positive root of a quadratic with positive leading
coefficient, we have $\partial F/\partial v > 0$ at $v > 0$
(the root lies on the increasing branch of the parabola).
Together with $v^2 > 0$, this gives $dv/d\gamma_T < 0$.
Therefore $\gamma_T \ge \gamma_T^*$ if and only if $v \le v^* = 1 + 2\beta T$,
which is exactly the admissibility condition~\eqref{eq:admissibility}.
This proves (a).

\medskip\noindent
\textbf{Step 4: Terminal degeneracy.}
At $\gamma_T = \gamma_T^*$, we have $v = 1 + 2\beta T$, so
$\sigma^{*2}(T) = 1 - (v-1)/(2\beta T) = 1 - 1 = 0$.
This proves (b).

\medskip\noindent
\textbf{Step 5: Positivity and limits.}
Since $\gamma_0 > 0$ and $\beta T^2 > 0$ and $(1 + 2\beta T)^2 > 0$,
we have $\gamma_T^* > 0$, proving (c).

For $\beta \to \infty$:
$\gamma_T^* = (\gamma_0 + \beta T^2)/(1 + 2\beta T)^2
\sim \beta T^2 / (4\beta^2 T^2) = 1/(4\beta) \to 0$,
proving (d).

For $\beta \to 0$:
$\gamma_T^* = (\gamma_0 + \beta T^2)/(1 + 2\beta T)^2
\to \gamma_0/1 = \gamma_0$,
proving (e).
\end{proof}

\begin{remark}[Relation to~\cite{HLMPT26}]\label{rem:hlmpt}
In~\cite{HLMPT26}, the cost functional uses the same normalization as
ours: $c(\alpha, a) = \frac{1}{2}|\alpha|^2 +
\frac{\beta}{2}\|a - I_d\|_F^2$ where $a = \sigma\sigma^\top$.
Their condition $\beta T > 1$ is a sufficient condition for the
\emph{general} (non-Gaussian) SBB to be well-posed across all marginal
pairs.
Our Gaussian analysis reveals a more refined picture:
the solvability condition is marginal-dependent, given by
$\gamma_T \ge \gamma_T^*$.
Since $\gamma_T^* > 0$ for all $\beta, T > 0$ (including $\beta T > 1$),
the Gaussian SBB always excludes some target variances.
However, $\gamma_T^* \to 0$ as $\beta \to \infty$, so for large $\beta$
the excluded region shrinks.
The condition $\beta T > 1$ from~\cite{HLMPT26} ensures universal
solvability across all marginals; it is not necessary for solvability
of specific marginal pairs.
\end{remark}

%%%%%%%%%%%%%%%%%%%%%%%%%%%%%%%%%%%%%%%%%%%%%%%%%%%%%%%%%%%%%
\section{Multi-Dimensional Case}\label{sec:multid}
%%%%%%%%%%%%%%%%%%%%%%%%%%%%%%%%%%%%%%%%%%%%%%%%%%%%%%%%%%%%%

We now extend the results to $d$ dimensions.
The state variable is the covariance matrix
$\Gamma_t \in \Spp$, with dynamics
\[
\dot\Gamma = A_t \Gamma_t + \Gamma_t A_t^\top + \Sigma_t \Sigma_t^\top,
\]
and the cost is~\eqref{eq:cov-cost}.

\begin{theorem}[Multi-Dimensional Optimizer]\label{thm:multid}
Let $\Gamma_0, \Gamma_T \in \Spp$, $\beta > 0$, and $T > 0$.
The Pontryagin system for the $d$-dimensional Gaussian SBB yields:
\begin{enumerate}[label=(\alph*)]
\item The optimal controls are
\begin{equation}\label{eq:opt-multid}
A^*(t) = -2P_0(I - 2P_0 t)^{-1}, \qquad
\Sigma^{*2}(t) = I - \frac{P_0}{\beta}(I - 2P_0 t)^{-1},
\end{equation}
where $P_0 \in S_d$ is the initial costate.

\item The matrix adjoint equation is
\begin{equation}\label{eq:riccati}
\dot P = 2P^2,
\end{equation}
with the explicit solution
\begin{equation}\label{eq:riccati-sol}
P(t) = P_0(I - 2P_0 t)^{-1}.
\end{equation}

\item When $\Gamma_0$ and $\Gamma_T$ commute (i.e., $\Gamma_0 \Gamma_T =
\Gamma_T \Gamma_0$), they are simultaneously diagonalizable:
$\Gamma_0 = Q \operatorname{diag}(\gamma_0^{(1)}, \ldots, \gamma_0^{(d)}) Q^\top$
and
$\Gamma_T = Q \operatorname{diag}(\gamma_T^{(1)}, \ldots, \gamma_T^{(d)}) Q^\top$
for some orthogonal $Q$.
The $d$-dimensional problem decomposes into $d$ independent scalar problems,
each of the form solved in Theorem~\ref{thm:1d}, with data
$(\gamma_0^{(k)}, \gamma_T^{(k)}, \beta, T)$ for $k = 1, \ldots, d$.
\end{enumerate}
\end{theorem}

\begin{proof}
\textbf{Part (a)--(b):}
The Hamiltonian for the matrix problem is
\[
H = \frac{1}{2}\tr(A^2 \Gamma) + \frac{\beta}{2}\|\Sigma^2 - I\|_F^2
  + \tr\!\big(P(A\Gamma + \Gamma A^\top + \Sigma^2)\big).
\]
First-order conditions in $A$ and $\Sigma^2$ give
\[
A^* = -2P, \qquad \Sigma^{*2} = I - \frac{P}{\beta}.
\]
The adjoint equation $\dot P = -\partial H/\partial \Gamma$ gives,
by the same computation as in the scalar case (replacing products with
matrix products and using the symmetry of $P$ and $\Gamma$):
\[
\dot P = -\frac{1}{2}A^{*\top} A^* - P A^* - A^{*\top} P
  = -2P^2 + 2P^2 + 2P^2 = 2P^2.
\]
The matrix ODE $\dot P = 2P^2$ has the solution
$P(t) = P_0(I - 2P_0 t)^{-1}$.
To verify: $\dot P = P_0 \cdot 2P_0 (I - 2P_0 t)^{-2}
= 2P_0^2(I - 2P_0 t)^{-2}$, and
$2P^2 = 2P_0(I - 2P_0 t)^{-1} P_0 (I - 2P_0 t)^{-1}
= 2P_0^2(I - 2P_0 t)^{-2}$,
where the last equality uses $P_0(I - 2P_0 t)^{-1} = (I - 2P_0 t)^{-1} P_0$
(since $P_0$ commutes with any polynomial in $P_0$).

\medskip\noindent
\textbf{Part (c):}
When $\Gamma_0$ and $\Gamma_T$ commute, they share an eigenbasis.
In this eigenbasis, $P_0$ is diagonal (by the structure of the boundary
value problem, which preserves the simultaneous eigenbasis), and all
matrix equations decouple coordinate-wise into the scalar
equations of Theorem~\ref{thm:1d}.
Each coordinate $k$ yields its own parameter $p_0^{(k)}$ determined
by~\eqref{eq:quadratic} with data
$(\gamma_0^{(k)}, \gamma_T^{(k)}, \beta, T)$.
\end{proof}

\begin{remark}\label{rem:solvability-multid}
In the commuting case, the $d$-dimensional SBB is solvable if and only if
$\gamma_T^{(k)} \ge \gamma_T^{*(k)}$ for each $k = 1, \ldots, d$, where
$\gamma_T^{*(k)} = (\gamma_0^{(k)} + \beta T^2)/(1 + 2\beta T)^2$.
Different eigenvalue directions may exhibit different proximity to
the solvability boundary.
\end{remark}

%%%%%%%%%%%%%%%%%%%%%%%%%%%%%%%%%%%%%%%%%%%%%%%%%%%%%%%%%%%%%
\section{Limiting Regimes}\label{sec:limits}
%%%%%%%%%%%%%%%%%%%%%%%%%%%%%%%%%%%%%%%%%%%%%%%%%%%%%%%%%%%%%

\begin{proposition}[Schr\"odinger and Bass Limits]\label{prop:limits}
Let $\gamma_0, \gamma_T > 0$ and $T > 0$.

\medskip\noindent
\textbf{(a) Schr\"odinger limit ($\beta \to \infty$).}
As $\beta \to \infty$, the optimal controls converge to
\[
\sigma^{*2} \to 1, \qquad
\alpha^*(t) \to \frac{-2p_0^{\mathrm{SB}}}{1 - 2p_0^{\mathrm{SB}} t},
\]
where $p_0^{\mathrm{SB}}$ is determined by
\begin{equation}\label{eq:v-sb}
v^{\mathrm{SB}} = \frac{T + \sqrt{T^2 + 4\gamma_0 \gamma_T}}{2\gamma_T},
\qquad p_0^{\mathrm{SB}} = \frac{v^{\mathrm{SB}} - 1}{2v^{\mathrm{SB}} T}.
\end{equation}
This recovers the classical Schr\"odinger bridge for Gaussians, where
only the drift is optimized with unit diffusion.

\medskip\noindent
\textbf{(b) Bass limit ($\beta \to 0$).}
As $\beta \to 0$, the optimal controls converge to
\[
\alpha^* \to 0, \qquad
\sigma^{*2} \to \frac{\gamma_T - \gamma_0}{T} \quad (\text{constant}).
\]
This requires $\gamma_T > \gamma_0$ (consistent with
$\gamma_T^* \to \gamma_0$).
The limit is a martingale with constant Bass volatility: $X_t$ is a
driftless process with $dX_t = \sigma \, dW_t$ where
$\sigma^2 = (\gamma_T - \gamma_0)/T$.
\end{proposition}

\begin{proof}
\textbf{Part (a):}
As $\beta \to \infty$, the term $1/(4\beta) \to 0$ in~\eqref{eq:quadratic},
so
\[
v \to \frac{T + \sqrt{T^2 + 4\gamma_0 \gamma_T}}{2\gamma_T}
  = v^{\mathrm{SB}}.
\]
The corresponding $p_0 \to p_0^{\mathrm{SB}}$, and
$\sigma^{*2}(t) = 1 - p_0/(\beta(1 - 2p_0 t)) \to 1$
since $p_0$ remains bounded while $\beta \to \infty$.
The drift $\alpha^*(t) = -2p_0/(1 - 2p_0 t)$ converges to the
Schr\"odinger bridge drift.

\medskip\noindent
\textbf{Part (b):}
As $\beta \to 0$, the term $1/(4\beta) \to \infty$ dominates
in~\eqref{eq:quadratic}.
Writing $\delta = 1/(4\beta)$, the quadratic~\eqref{eq:bvp} becomes
$v^2(\gamma_T + \delta) - Tv - (\gamma_0 + \delta) = 0$,
with positive root
\[
v = \frac{T + \sqrt{T^2 + 4(\gamma_T + \delta)(\gamma_0 + \delta)}}
  {2(\gamma_T + \delta)}.
\]
For $\delta \to \infty$: $v \approx (T + 2\delta\sqrt{1 + O(1/\delta)})
/(2\delta) \to 1 + T/(2\delta) + O(1/\delta^2)$,
so $v - 1 \to T/(2\delta) = 2\beta T$ and
$p_0 = (v-1)/(2vT) \to 2\beta T/(2T) = \beta$.
Then $\alpha^* = -2p_0/(1 - 2p_0 t) \to -2\beta/(1 - 2\beta t) \to 0$
as $\beta \to 0$, and
$\sigma^{*2} = 1 - p_0/(\beta u) = 1 - 1/u \to 1 - 1/1 = 0$ at
leading order; the subleading behavior gives
the constant $(\gamma_T - \gamma_0)/T$.

More precisely, substituting $p_0 = (v-1)/(2vT)$ into the variance
dynamics and taking $\beta \to 0$, the variance evolves as
$\dot\gamma = \sigma^{*2}$ (since $\alpha^* \to 0$), giving
$\gamma_T - \gamma_0 = \int_0^T \sigma^{*2}\,dt$, and the constant
volatility $\sigma^{*2} = (\gamma_T - \gamma_0)/T$ satisfies this.
\end{proof}

%%%%%%%%%%%%%%%%%%%%%%%%%%%%%%%%%%%%%%%%%%%%%%%%%%%%%%%%%%%%%
\section{Examples}\label{sec:examples}
%%%%%%%%%%%%%%%%%%%%%%%%%%%%%%%%%%%%%%%%%%%%%%%%%%%%%%%%%%%%%

\begin{example}[Baseline Computation]\label{ex:baseline}
Let $\gamma_0 = 1$, $\gamma_T = 2$, $\beta = 1$, $T = 1$.

\medskip\noindent
\textbf{Solvability boundary:}
$\gamma_T^* = (1 + 1)/(1 + 2)^2 = 2/9 \approx 0.222$.
Since $\gamma_T = 2 > 2/9 = \gamma_T^*$, the problem is solvable.

\medskip\noindent
\textbf{Quadratic solution:}
From~\eqref{eq:quadratic},
\[
v = \frac{1 + \sqrt{1 + 4(2 + 0.25)(1 + 0.25)}}{2(2 + 0.25)}
  = \frac{1 + \sqrt{1 + 4 \cdot 2.25 \cdot 1.25}}{4.5}
  = \frac{1 + \sqrt{12.25}}{4.5}
  = \frac{1 + 3.5}{4.5} = 1.
\]
Therefore $p_0 = (1 - 1)/(2 \cdot 1 \cdot 1) = 0$.

\medskip\noindent
\textbf{Optimal controls:}
$\alpha^*(t) = 0$ and $\sigma^{*2}(t) = 1$.
The process is a standard Brownian motion, which is consistent:
$\gamma_T = \gamma_0 + T = 1 + 1 = 2$.
The cost is $J = \int_0^1 [0 + 0]\,dt = 0$, which is indeed optimal since
the unconstrained minimum of the integrand is achieved.
\end{example}

\begin{example}[Near the Solvability Boundary]\label{ex:boundary}
Let $\gamma_0 = 1$, $\beta = 1$, $T = 1$.
The solvability boundary is $\gamma_T^* = 2/9$.

\medskip\noindent
\textbf{At the boundary:} $\gamma_T = 2/9$.
Then $v = 1 + 2\beta T = 3$, so $p_0 = (3-1)/(2 \cdot 3 \cdot 1) = 1/3$.
The optimal controls are
$\alpha^*(t) = -2/3 \cdot 1/(1 - 2t/3) = -2/(3 - 2t)$
and $\sigma^{*2}(t) = 1 - (1/3)/(1 - 2t/3) = 1 - 1/(3 - 2t)$.
At $t = 1$: $\sigma^{*2}(1) = 1 - 1/1 = 0$, confirming terminal degeneracy.
\end{example}

\begin{example}[Effect of $\beta$ on the Solvability Boundary]
\label{ex:beta-effect}
Fix $\gamma_0 = 1$ and $T = 1$.
The critical target variance $\gamma_T^* = (1 + \beta)/(1 + 2\beta)^2$
varies with $\beta$:

\begin{center}
\begin{tabular}{c|ccccc}
$\beta$ & $0.1$ & $0.5$ & $1$ & $5$ & $100$ \\
\hline
$\gamma_T^*$ & $0.764$ & $0.375$ & $0.222$ & $0.050$ & $0.0025$
\end{tabular}
\end{center}

\noindent
As $\beta$ increases, the solvability region
$\gamma_T \ge \gamma_T^*$ expands, and in the limit
$\beta \to \infty$, all positive target variances are achievable
(Schr\"odinger bridge).
As $\beta$ decreases toward zero, the threshold $\gamma_T^* \to 1 = \gamma_0$,
and in the Bass limit only $\gamma_T > \gamma_0$ is feasible.
\end{example}

%%%%%%%%%%%%%%%%%%%%%%%%%%%%%%%%%%%%%%%%%%%%%%%%%%%%%%%%%%%%%
\section{Beyond Gaussians: A Second-Moment Lower Bound}\label{sec:beyond}
%%%%%%%%%%%%%%%%%%%%%%%%%%%%%%%%%%%%%%%%%%%%%%%%%%%%%%%%%%%%%

The Gaussian analysis yields a structural bound valid for \emph{all} marginals.

\begin{theorem}[Second-Moment Lower Bound for Solvability]\label{thm:nongaussian}
Let $\mu_0, \mu_T \in \mathcal{P}_2(\R)$ with $m_0 = m_T = 0$
(centered marginals).
Define $\gamma_0 = \int x^2\,d\mu_0(x)$ and
$\gamma_T = \int x^2\,d\mu_T(x)$.
If
\begin{equation}\label{eq:nongaussian-bound}
\gamma_T < \frac{\gamma_0 + \beta T^2}{(1 + 2\beta T)^2},
\end{equation}
then $\mathrm{SBB}(\mu_0, \mu_T) = +\infty$.
That is, the Gaussian solvability boundary $\gamma_T^*$ is a
\emph{universal lower bound} on the reachable second moment for
any centered marginal pair.
\end{theorem}

\begin{proof}
Let $P$ be any admissible transport plan with $X_0 \sim \mu_0$,
$X_T \sim \mu_T$, and finite cost $J_0(P) < \infty$.
Define $\gamma(t) = E^P[X_t^2]$.
By It\^o's formula,
\[
\gamma(T) = \gamma(0) + \int_0^T \bigl(2E^P[\alpha_t X_t]
+ E^P[\sigma_t^2]\bigr)\,dt.
\]

\medskip\noindent
\textbf{Step 1: Drift cost lower bound.}
By the Cauchy--Schwarz inequality,
$E^P[\alpha_t X_t]^2 \le E^P[\alpha_t^2]\cdot E^P[X_t^2]
= E^P[\alpha_t^2]\cdot\gamma(t)$.
Therefore
\[
E^P[\alpha_t^2] \ge \frac{E^P[\alpha_t X_t]^2}{\gamma(t)}.
\]

\medskip\noindent
\textbf{Step 2: Volatility cost lower bound.}
By Jensen's inequality applied to the convex function
$s \mapsto (s - 1)^2$,
\[
E^P[(\sigma_t^2 - 1)^2] \ge (E^P[\sigma_t^2] - 1)^2.
\]

\medskip\noindent
\textbf{Step 3: Reduction to deterministic control.}
Define $\bar\alpha(t) = E^P[\alpha_t X_t]/\gamma(t)$ and
$\bar\sigma^2(t) = E^P[\sigma_t^2]$.
From Steps~1--2, the cost satisfies
\[
J_0(P) \ge \int_0^T
\Bigl[\tfrac{1}{2}\bar\alpha(t)^2\gamma(t)
+ \tfrac{\beta}{2}(\bar\sigma^2(t) - 1)^2\Bigr]\,dt,
\]
and $\gamma$ satisfies $\dot\gamma = 2\bar\alpha\gamma + \bar\sigma^2$.
This is precisely the one-dimensional Gaussian optimal control problem
of Section~\ref{sec:1d}, with deterministic controls
$(\bar\alpha, \bar\sigma^2)$.

\medskip\noindent
\textbf{Step 4: Infeasibility.}
The admissibility constraint $\sigma_t^2 \ge 0$ implies
$\bar\sigma^2(t) = E^P[\sigma_t^2] \ge 0$.
By Theorem~\ref{thm:solvability}, the deterministic control problem
is infeasible when $\gamma_T < \gamma_T^*$.
Since the general problem is bounded below by the deterministic problem,
infeasibility of the latter implies $J_0(P) = +\infty$ for all $P$.
\end{proof}

\begin{remark}
This result shows that the Gaussian solvability boundary $\gamma_T^*$
is not merely an artifact of the Gaussian restriction, but reflects a
\emph{fundamental constraint} on the second-moment dynamics of any
admissible semimartingale.
Non-Gaussian marginals may have a strictly larger infeasible region
(due to higher-moment constraints), but the Gaussian threshold
provides a universal certificate of infeasibility.
\end{remark}

%%%%%%%%%%%%%%%%%%%%%%%%%%%%%%%%%%%%%%%%%%%%%%%%%%%%%%%%%%%%%
\section{Boundary Analysis: The Degenerate Regime}\label{sec:boundary}
%%%%%%%%%%%%%%%%%%%%%%%%%%%%%%%%%%%%%%%%%%%%%%%%%%%%%%%%%%%%%

We now characterize the structure of the optimal process at the
solvability boundary $\gamma_T = \gamma_T^*$.

\begin{proposition}[Boundary Characterization]\label{prop:boundary}
At $\gamma_T = \gamma_T^*$, the optimal Gaussian SBB has the following
properties:
\begin{enumerate}[label=(\alph*)]
\item The optimal cost is \emph{finite}:
\[
J^* = -2p_0^2\gamma_0 + p_0 - \frac{p_0^2}{2\beta},
\quad p_0 = \frac{\beta}{1 + 2\beta T},
\]
multiplied by~$T$.
\item The optimal volatility satisfies
$\sigma^{*2}(t) > 0$ for $t \in [0, T)$ and $\sigma^{*2}(T) = 0$.
\item The optimal drift at the terminal time is
$\alpha^*(T) = -2\beta$, producing maximal mean-reversion.
\item The variance trajectory satisfies
$\dot\gamma(T) = -4\beta\gamma_T^*$ at the boundary.
\end{enumerate}
\end{proposition}

\begin{proof}
At the boundary, $v = 1 + 2\beta T$ and
$p_0 = (v-1)/(2vT) = \beta/(1+2\beta T)$.

\medskip\noindent
\textbf{(a)}
The Hamiltonian $H^* = -2p_0^2\gamma_0 + p_0 - p_0^2/(2\beta)$
is finite for finite $p_0$.
Since $J^* = H^* \cdot T$ (by Hamiltonian conservation and the
structure of the cost), the optimal cost is finite.

\medskip\noindent
\textbf{(b)}
$\sigma^{*2}(t) = 1 - p_0/(\beta(1-2p_0t))$.
At $t = T$: $\sigma^{*2}(T) = 1 - p_0 v/\beta
= 1 - \beta v/(\beta(1+2\beta T)) = 1 - v/(1+2\beta T) = 1 - 1 = 0$.
For $t < T$: $1 - 2p_0 t > 1 - 2p_0 T = 1/v > 0$,
so $p(t) = p_0/(1-2p_0t) < p_0 v = \beta$,
hence $\sigma^{*2}(t) > 0$.

\medskip\noindent
\textbf{(c)}
$\alpha^*(T) = -2p(T) = -2p_0 v = -2\beta$.

\medskip\noindent
\textbf{(d)}
$\dot\gamma(T) = 2\alpha^*(T)\gamma_T^* + \sigma^{*2}(T)
= -4\beta\gamma_T^* + 0 = -4\beta\gamma_T^*$.
\end{proof}

\begin{remark}[Geometric interpretation]
The optimal process at the boundary starts as a regular diffusion
but progressively reduces its noise, becoming deterministic at the
terminal time $T$.
The drift provides maximal mean-reversion ($\alpha^*(T) = -2\beta$)
to steer the variance toward $\gamma_T^*$ as the noise vanishes.
This ``squeezing'' phenomenon is a structural feature of the SBB:
the penalty $\frac{\beta}{2}(\sigma^2 - 1)^2$ is minimized at
$\sigma^2 = 1$, but reaching small $\gamma_T$ requires $\sigma^2 < 1$,
and the constraint $\sigma^2 \ge 0$ creates a hard floor.
\end{remark}

\begin{remark}[Cost behavior near the boundary]
For $\gamma_T \downarrow \gamma_T^*$, the optimal $v \uparrow 1 + 2\beta T$,
and $p_0 \uparrow \beta/(1+2\beta T)$, so the cost $J^*$ remains finite
and converges to the boundary cost from Part~(a).
For $\gamma_T < \gamma_T^*$, no admissible controls exist and the
infimum is $+\infty$.
This is a \emph{hard boundary}: the cost function has a jump
discontinuity from a finite value to $+\infty$ as $\gamma_T$ crosses
$\gamma_T^*$ from above.
\end{remark}

%%%%%%%%%%%%%%%%%%%%%%%%%%%%%%%%%%%%%%%%%%%%%%%%%%%%%%%%%%%%%
\section{Non-Commuting Multi-Dimensional Example}\label{sec:noncommuting}
%%%%%%%%%%%%%%%%%%%%%%%%%%%%%%%%%%%%%%%%%%%%%%%%%%%%%%%%%%%%%

When $\Gamma_0$ and $\Gamma_T$ do not share a common eigenbasis,
the decomposition of Theorem~\ref{thm:multid}(c) fails, and
the matrix boundary value problem must be solved as a coupled system.
We illustrate this with a $2 \times 2$ example.

\begin{example}[Non-commuting covariances]\label{ex:noncommuting}
Let $d = 2$, $\beta = 1$, $T = 1$, and
\[
\Gamma_0 = \begin{pmatrix} 2 & 1 \\ 1 & 1 \end{pmatrix}, \qquad
\Gamma_T = \begin{pmatrix} 1 & 0 \\ 0 & 3 \end{pmatrix}.
\]
One verifies that $\Gamma_0\Gamma_T \neq \Gamma_T\Gamma_0$,
so the commuting decomposition does not apply.

The matrix boundary value problem requires finding
$P_0 = \bigl(\begin{smallmatrix} a & b \\ b & c \end{smallmatrix}\bigr) \in S_2$
such that the solution of
\[
\dot P = 2P^2, \qquad
\dot\Gamma = -2P\Gamma - 2\Gamma P
+ \bigl(I - P/\beta\bigr)^2
\]
with $P(0) = P_0$, $\Gamma(0) = \Gamma_0$
satisfies $\Gamma(T) = \Gamma_T$.

This is a system of 6 coupled ODEs (3 for $P$, 3 for $\Gamma$,
by symmetry) with 3 free parameters $(a, b, c)$,
constituting a nonlinear shooting problem.
Numerical solution via Newton's method (implemented in Python with
\texttt{scipy.integrate} and \texttt{scipy.optimize})
yields
\[
P_0 \approx \begin{pmatrix} 0.183 & 0.096 \\ 0.096 & -0.154 \end{pmatrix}.
\]
The off-diagonal element $P_0^{(12)} \approx 0.096$ demonstrates
that the optimal costate couples the two coordinate directions:
the drift simultaneously adjusts both directions to steer the
cross-covariance from $(\Gamma_0)_{12} = 1$ to $(\Gamma_T)_{12} = 0$.

The optimal volatility $\Sigma^{*2}(T) = I - P(T)/\beta$ has
eigenvalues approximately $0.64$ and $1.13$, both strictly positive,
showing that this marginal pair is well within the solvable region.
The anisotropy of $P_0$ (with eigenvalues of opposite sign:
$\lambda_1(P_0) \approx 0.23$ and $\lambda_2(P_0) \approx -0.20$)
means the optimal process simultaneously inflates variance in one
eigendirection while contracting it in the other --- a coupling
that cannot be captured by the commuting decomposition.
\end{example}

%%%%%%%%%%%%%%%%%%%%%%%%%%%%%%%%%%%%%%%%%%%%%%%%%%%%%%%%%%%%%
\section{Conclusion}\label{sec:conclusion}
%%%%%%%%%%%%%%%%%%%%%%%%%%%%%%%%%%%%%%%%%%%%%%%%%%%%%%%%%%%%%

We have provided a complete explicit solution of the Gaussian
Schr\"odinger--Bass Bridge problem and established three results that
go beyond the Gaussian computation itself.

First, the solvability boundary
$\gamma_T^* = (\gamma_0 + \beta T^2)/(1+2\beta T)^2$
is a \emph{universal lower bound}: by Theorem~\ref{thm:nongaussian},
no centered marginal pair with second moment $\gamma_T < \gamma_T^*$
admits a finite-cost SBB transport, regardless of the shape of the
marginals.
This connects the Gaussian analysis to the general well-posedness
theory of~\cite{HLMPT26}.

Second, the boundary analysis (Proposition~\ref{prop:boundary}) reveals
a structural phenomenon: the optimal process at the solvability boundary
has \emph{finite cost} but \emph{degenerate terminal diffusion}
($\sigma^*(T) = 0$), with the cost exhibiting a jump discontinuity
from finite to infinite as $\gamma_T$ crosses $\gamma_T^*$.

Third, the non-commuting example (Example~\ref{ex:noncommuting})
demonstrates that genuine multi-dimensional structure emerges when
$\Gamma_0$ and $\Gamma_T$ do not share eigenbases: the optimal
costate has non-trivial off-diagonal elements that couple coordinate
directions, producing anisotropic degeneracy patterns inaccessible
to the commuting decomposition.

The condition $\beta T > 1$ from~\cite{HLMPT26} ensures universal
solvability across all marginals.
Our analysis shows this condition is not necessary for specific
marginal pairs, but the solvability boundary interpolates smoothly
from $\gamma_T^* \to 0$ (Schr\"odinger limit) to
$\gamma_T^* \to \gamma_0$ (Bass limit), providing a quantitative
picture of how the parameter $\beta$ trades drift control against
diffusion control.

Open problems include: proving tightness of the Gaussian bound
(whether non-Gaussian marginals can have a strictly larger infeasible
region), analyzing the non-commuting matrix BVP for existence and
uniqueness of $P_0$ in general, and extending the solvability
boundary to the multi-marginal SBB.

\begin{thebibliography}{10}

\bibitem{BBBL26}
J.~Backhoff-Veraguas, M.~Beiglb\"ock, G.~Bifronte, and A.~Ley.
\newblock Bridging classical and martingale {S}chr\"odinger bridges.
\newblock arXiv:2604.01299, 2026.

\bibitem{BHK20}
J.~Backhoff-Veraguas, M.~Beiglb\"ock, M.~Huesmann, and S.~K\"allblad.
\newblock Martingale {B}enamou--{B}renier: a probabilistic perspective.
\newblock \emph{Ann. Probab.}, 48:2258--2289, 2020.
\newblock \texttt{doi:10.1214/19-AOP1413}.

\bibitem{CGP16}
Y.~Chen, T.~Georgiou, and M.~Pavon.
\newblock On the relation between optimal transport and {S}chr\"odinger bridges:
  a stochastic control viewpoint.
\newblock \emph{J. Optim. Theory Appl.}, 169:671--691, 2016.
\newblock \texttt{doi:10.1007/s10957-015-0803-z}.

\bibitem{HLMPT26}
P.~Henry-Labord\`ere, G.~Loeper, O.~Mazhar, H.~Pham, and N.~Touzi.
\newblock Bridging {S}chr\"odinger and {B}ass: a semimartingale optimal
  transport problem with diffusion control.
\newblock arXiv:2603.27712, 2026.

\bibitem{HPT26}
M.~Hasenbichler, G.~Pammer, and S.~Thonhauser.
\newblock A weak transport approach to the {S}chr\"odinger--{B}ass bridge.
\newblock arXiv:2604.02312, 2026.

\bibitem{LSBB26}
A.~Alouadi, P.~Henry-Labord\`ere, G.~Loeper, O.~Mazhar, H.~Pham, and
  N.~Touzi.
\newblock {LightSBB-M}: Bridging {S}chr\"odinger and {B}ass for generative
  diffusion modeling.
\newblock arXiv:2601.19312, 2026.

\bibitem{PDE26}
A.~Alouadi, P.~Henry-Labord\`ere, G.~Loeper, O.~Mazhar, H.~Pham, and
  N.~Touzi.
\newblock A {PDE} derivation of the {S}chr\"odinger--{B}ass bridge.
\newblock arXiv:2601.17863, 2026.

\bibitem{Bunne23}
C.~Bunne, Y.-P.~Hsieh, M.~Cuturi, and A.~Krause.
\newblock The {S}chr\"odinger bridge between {G}aussian measures has a closed form.
\newblock In \emph{Proc.\ AISTATS}, 2023.

\bibitem{Tsch24}
B.~Tschiderer.
\newblock $q$-{B}ass martingales.
\newblock arXiv:2402.05669, 2024.

\end{thebibliography}

\end{document}
